\documentclass{article}
\usepackage{graphicx} % Extended graphics inclusions
\usepackage{float}
\usepackage{url} % For \url{}
\usepackage{../config/atxy} % For front cover
\usepackage{amsfonts} % Needed for some fonts
\usepackage[usenames]{color} % Needed for colored R input/output
\usepackage{pdfcolmk} % Correct some problems with the color stack


\title{Nonparametric statistics}
\author{Palmeira, L. \and Lobry, J.R.}

\usepackage{/Library/Frameworks/R.framework/Resources/share/texmf/Sweave}
\begin{document}
%
% To change the R input/output style:
%
\definecolor{Soutput}{rgb}{0,0,0.56}
\definecolor{Sinput}{rgb}{0.56,0,0}
\DefineVerbatimEnvironment{Sinput}{Verbatim}
{formatcom={\color{Sinput}},fontsize=\footnotesize, baselinestretch=0.75}
\DefineVerbatimEnvironment{Soutput}{Verbatim}
{formatcom={\color{Soutput}},fontsize=\footnotesize, baselinestretch=0.75}
%
% Rlogo
%
\newcommand{\Rlogo}{\protect\includegraphics[height=1.8ex,keepaspectratio]{../figs/Rlogo.pdf}}
%
% Shortcut for seqinR:
%
\newcommand{\seqinr}{\texttt{seqin\bf{R}}}
\newcommand{\Seqinr}{\texttt{Seqin\bf{R}}}
\fvset{fontsize= \scriptsize}
%
% R output options and libraries to be loaded.
%
%
%  Sweave Options
%
% Put all figures in the fig folder and start the name with current file name.
% Do not produce EPS files
%


\maketitle
% BEGIN - DO NOT REMOVE THIS LINE

%
% Prochlorococcus marinus example
%

\section{Introduction}

Nonparametric statistical methods were initially developped to study
variables for which little or nothing is known concerning their
distribution. This makes them particularly suitable for statistical
analysis of biological sequences, in particular for the study of over-
and under-representation of $k$-letter words (\textit{cf} section 
page \ref{dinu}).

\section{Elementary nonparametric statistics}

\subsection{Introduction}

Those rank statistics are those that were available under the ANALSEQ
software \cite{analseq, acnuc1984}. We consider here a sequence
of booleans, for instance:

\begin{Schunk}
\begin{Sinput}
 (x <- rep(c(T, F), 10))
\end{Sinput}
\begin{Soutput}
 [1]  TRUE FALSE  TRUE FALSE  TRUE FALSE  TRUE FALSE  TRUE FALSE  TRUE FALSE
[13]  TRUE FALSE  TRUE FALSE  TRUE FALSE  TRUE FALSE
\end{Soutput}
\end{Schunk}

We note $\mathrm{N}$ the total number of elements in the vector:

\begin{Schunk}
\begin{Sinput}
 (N <- length(x))
\end{Sinput}
\begin{Soutput}
[1] 20
\end{Soutput}
\end{Schunk}

We note $\mathrm{M}$ the total number of TRUE elements in the vector:

\begin{Schunk}
\begin{Sinput}
 (M <- sum(x))
\end{Sinput}
\begin{Soutput}
[1] 10
\end{Soutput}
\end{Schunk}

We note $\omega$ the ranks of TRUE elements:

\begin{Schunk}
\begin{Sinput}
 (omega <- which(x))
\end{Sinput}
\begin{Soutput}
 [1]  1  3  5  7  9 11 13 15 17 19
\end{Soutput}
\end{Schunk}

With one exception, the statistics names are the same as in the ANALSEQ software.

\subsection{Rank sum}

The statistic $\mathrm{SR}$ is the sum of the ranks of TRUE elements.

$$
\mathrm{SR} = \sum_{j \in \omega}{j}
$$

\begin{verbatim}
**-***-----------     ==> SR low  (18)
---------***--***     ==> SR high (81)
\end{verbatim}

\begin{eqnarray*}
\mathrm{E(SR)} & = & \mathrm{\frac{M(N + 1)}{2}} \\
\mathrm{V(SR)} & = & \mathrm{\frac{M(N + 1)(N - M)}{12}}
\end{eqnarray*}


\begin{Schunk}
\begin{Sinput}
 SR <- function(bool, N = length(bool), M = sum(bool)) {
     stopifnot(is.logical(bool))
     SR <- sum(seq_len(N)[bool])
     E <- M * (N + 1)/2
     V <- M * (N + 1) * (N - M)/12
     return(list(SR = SR, stat = (SR - E)/sqrt(V)))
 }
 SR(s2c("**-***-----------") == "*")
\end{Sinput}
\begin{Soutput}
$SR
[1] 18

$stat
[1] -2.84605
\end{Soutput}
\begin{Sinput}
 SR(s2c("---------***--***") == "*")
\end{Sinput}
\begin{Soutput}
$SR
[1] 81

$stat
[1] 2.713602
\end{Soutput}
\end{Schunk}

Here is a way to obtain the same result using the standard \Rlogo{} \texttt{wilcox.test()}
function to make a Wilcoxon's rank sum test \cite{WilcoxonF1945}:

\begin{Schunk}
\begin{Sinput}
 SRh <- s2c("---------***--***") == "*"
 x <- seq_len(length(SRh))
 x[!SRh] <- -1 * x[!SRh]
 wilcox.test(x)$statistic
\end{Sinput}
\begin{Soutput}
 V 
81 
\end{Soutput}
\end{Schunk}

The probabilities for all possibe outcomes for the rank sums are given by
\texttt{dwilcox()} but note the $\frac{M(M+1)}{2}$ shift:

\begin{Schunk}
\begin{Sinput}
 m <- sum(SRh)
 n <- length(SRh) - m
 pdf <- dwilcox(x = 0:(n * m), m = m, n = n)
 plot(x = 0:(m * n) + m * (m + 1)/2, y = pdf, xlab = "Possible rank sums", 
     ylab = "Density", main = paste("---------***--*** : N =", 
         length(SRh), "M =", sum(SRh)), pch = 19)
 points(SR(SRh)$SR, dwilcox(x = SR(SRh)$SR - m * (m + 1)/2, 
     m = m, n = n), col = "red", pch = 19)
 arrows(x0 = SR(SRh)$SR, y0 = 0.01, x1 = SR(SRh)$SR, y1 = 0.0015, 
     length = 0.1)
 text(SR(SRh)$SR, 0.01, "Observed\nvalue", pos = 3)
\end{Sinput}
\end{Schunk}
\includegraphics{../figs/nonparastats-demodwilcox}

\subsection{Rank variance}

This statistics is the variance of ranks:

$$
\mathrm{VR} = \sum_{j \in \omega}{(j - \mathrm{\frac{N+1}{2}})^2}
$$

\begin{verbatim}
------****-------    ==> VR low  (6)
***----------****    ==> VR high (323)
\end{verbatim}

\begin{eqnarray*}
\mathrm{E(VR)} & = & \mathrm{\frac{M(N + 1)(N-1)}{12}} \\
\mathrm{V(VR)} & = & \mathrm{\frac{M(N - M)(N + 1)(N + 2)(N - 2)}{180}}
\end{eqnarray*}

\begin{Schunk}
\begin{Sinput}
 VR <- function(bool, N = length(bool), M = sum(bool)) {
     stopifnot(is.logical(bool))
     VR <- sum((seq_len(N)[bool] - (N + 1)/2)^2)
     E <- (M * (N + 1) * (N - 1))/12
     V <- (M * (N - M) * (N + 1) * (N + 2) * (N - 2))/180
     return(list(VR = VR, stat = (VR - E)/sqrt(V)))
 }
 VR(s2c("------****-------") == "*")
\end{Sinput}
\begin{Soutput}
$VR
[1] 6

$stat
[1] -2.337860
\end{Soutput}
\begin{Sinput}
 VR(s2c("***----------****") == "*")
\end{Sinput}
\begin{Soutput}
$VR
[1] 323

$stat
[1] 3.470246
\end{Soutput}
\end{Schunk}

We can use simulations to have an idea of the probability density function
of the rank variance, for instance:

\begin{Schunk}
\begin{Sinput}
 VRh <- s2c("***----------****") == "*"
 simVR <- replicate(5000, VR(sample(VRh))$VR)
 hist(simVR, col = grey(0.7), main = paste("***----------**** : N =", 
     length(VRh), "M =", sum(VRh)), xlab = "Possible rank variances", 
     proba = TRUE)
 lines(density(simVR), lwd = 2)
 arrows(VR(VRh)$VR, 0.004, VR(VRh)$VR, 0, le = 0.1)
\end{Sinput}
\end{Schunk}
\includegraphics{../figs/nonparastats-simVR}

\subsection{Clustering around the observed centre}

Let note $\mathrm{C(\omega)}$ the observed centre:

$$
\mathrm{C(\omega)} = \left\{ \begin{array}{rl}
 \omega\left(\mathrm{\frac{M+1}{2}}\right) &\mbox{ if $\mathrm{M}$ is odd} \\
 \omega\left(\mathrm{\frac{M}{2} + 1}\right) &\mbox{ if $\mathrm{M}$ is even}
       \end{array} \right.
$$

The statistic $\mathrm{CC}$\footnote{
the original notation was GC in the ANALSEQ software, we use CC instead to
avoid a collision with the \texttt{GC()} function to compute the G+C content.
} is the dispersion around $\mathrm{C(\omega)}$ is defined by:

$$
\mathrm{CC} = \sum_{j \in \omega}{\left| j - \mathrm{C(\omega)} \right|}
$$

\begin{verbatim}
---*****---------    ==> CC low  (6)
***-------***----    ==> CC high (30)
\end{verbatim}

Noting $\left[x\right]$ the floor of $x$, we have:

$$
\mathrm{E(CC)} = \mathrm{\frac{(N+1)\left[\frac{M}{2}\right]\left[\frac{M+1}{2}\right]}{M+1}}
$$

and

$$
\mathrm{V(CC)} = \left\{ \begin{array}{rl}
 \mathrm{\frac{(M - 1)(M + 3)(N + 1)(N - M)}{48(M + 2)}} &\mbox{ if $\mathrm{M}$ is odd} \\
 \mathrm{\frac{M(N + 1)(N - M)(M^2 + 2*M + 4)}{48(M + 1)^2}} &\mbox{ if $\mathrm{M}$ is even}
       \end{array} \right.
$$

\begin{Schunk}
\begin{Sinput}
 CC <- function(bool, N = length(bool), M = sum(bool)) {
     stopifnot(is.logical(bool))
     C <- median(seq_len(N)[bool])
     GC <- sum(abs(seq_len(N)[bool] - C))
     E <- ((N + 1) * floor(M/2) * floor((M + 1)/2))/(M + 1)
     if (M%%2 == 1) 
         V <- ((M - 1) * (M + 3) * (N + 1) * (N - M))/(48 * 
             (M + 2))
     else V <- (M * (N + 1) * (N - M) * (M^2 + 2 * M + 4))/(48 * 
         (M + 1)^2)
     return(list(GC = GC, stat = (GC - E)/sqrt(V)))
 }
 CC(s2c("---*****---------") == "*")
\end{Sinput}
\begin{Soutput}
$GC
[1] 6

$stat
[1] -2.645751
\end{Soutput}
\begin{Sinput}
 CC(s2c("***-------***----") == "*")
\end{Sinput}
\begin{Soutput}
$GC
[1] 30

$stat
[1] 1.337987
\end{Soutput}
\end{Schunk}

\subsection{Number of runs}

The statistics $\mathrm{NS}$ is the number of runs in the sequence:

\begin{verbatim}
--***---***--***-    ==> NS low  (7)
-*-*-*-*-*-*-*-*-    ==> NS high (17)
\end{verbatim}

\begin{eqnarray*}
\mathrm{E(NS)} & = & \mathrm{\frac{2M(N - M)}{N} + 1} \\
\mathrm{V(NS)} & = & \mathrm{\frac{2M(N - M)(2M(N - M) - N)}{N^2(N - 1)}}
\end{eqnarray*}

\begin{Schunk}
\begin{Sinput}
 NS <- function(bool, N = length(bool), M = sum(bool)) {
     stopifnot(is.logical(bool))
     NS <- length(rle(bool)$lengths)
     DMNmM <- 2 * M * (N - M)
     E <- DMNmM/N + 1
     V <- (DMNmM * (DMNmM - N))/(N * N * (N - 1))
     return(list(NS = NS, stat = (NS - E)/sqrt(V)))
 }
 NS(s2c("--***---***--***-") == "*")
\end{Sinput}
\begin{Soutput}
$NS
[1] 7

$stat
[1] -1.242299
\end{Soutput}
\begin{Sinput}
 NS(s2c("-*-*-*-*-*-*-*-*-") == "*")
\end{Sinput}
\begin{Soutput}
$NS
[1] 17

$stat
[1] 3.786054
\end{Soutput}
\end{Schunk}

The same result can be obtained with the function \texttt{runs.test()} from
package \textbf{tseries} \cite{tseries} this way:

\begin{Schunk}
\begin{Sinput}
 library(tseries)
 NSh <- s2c("-*-*-*-*-*-*-*-*-") == "*"
 tseries::runs.test(as.factor(NSh))$statistic
\end{Sinput}
\begin{Soutput}
Standard Normal 
       3.786054 
\end{Soutput}
\end{Schunk}

\subsection{Multiple clusters}

The statistics $\mathrm{GM}$ is the variance of the length $n_k$ of FALSE runs (including
runs of length zero) between two TRUE. Let note:

\begin{itemize}
\item $n_k(\omega)$ the number of FALSE between $\omega(k-1)$ and $\omega(k)$ for $2 \le k \le \mathrm{M}$.
\item $n_1(\omega)$ the number of FALSE before $\omega(1)$.
\item $n_{\mathrm{M+1}}(\omega)$ the number of FALSE after $\omega(\mathrm{M})$.
\end{itemize}

$$
\mathrm{GM} = \frac{1}{\mathrm{M}}\sum_{i = 1}^{\mathrm{M+1}}{\left( n_i(\omega) - \mathrm{\frac{N-M}{M+1}} \right)^2}
$$

\begin{verbatim}
-*-*-*-*-*-*-*-*-    ==> GM low  (0)
***------***-***-    ==> GM high (3.5)
\end{verbatim}

\begin{eqnarray*}
\mathrm{E(GM)} & = & \mathrm{\frac{(N + 1)(N - M)}{(M +1)(M + 2)}} \\
\mathrm{V(GM)} & = & \mathrm{\frac{4(N - M - 1)(N + 1)(N + 2)(N - M)}{M(M + 2)^2(M + 3)(M + 4)}}
\end{eqnarray*}

\begin{Schunk}
\begin{Sinput}
 GM <- function(bool, N = length(bool), M = sum(bool)) {
     stopifnot(is.logical(bool))
     XGM <- (N - M)/(M + 1)
     LS0 <- GM <- 0
     for (i in seq_len(N)) {
         if (bool[i]) {
             GM <- GM + (LS0 - XGM)^2
             LS0 <- 0
         }
         else {
             LS0 <- LS0 + 1
         }
     }
     GM <- (GM + (LS0 - XGM)^2)/M
     E <- ((N + 1) * (N - M))/((M + 1) * (M + 2))
     V <- (4 * (N - M - 1) * (N + 1) * (N + 2) * (N - M))/(M * 
         (M + 2)^2 * (M + 3) * (M + 4))
     return(list(GM = GM, stat = (GM - E)/sqrt(V)))
 }
 GM(s2c("-*-*-*-*-*-*-*-*-") == "*")
\end{Sinput}
\begin{Soutput}
$GM
[1] 0

$stat
[1] -1.863782
\end{Soutput}
\begin{Sinput}
 GM(s2c("***------***-***-") == "*")
\end{Sinput}
\begin{Soutput}
$GM
[1] 3.511111

$stat
[1] 3.279144
\end{Soutput}
\end{Schunk}

\section{Dinucleotides over- and under-representation}
\label{dinu}

\subsection{Introduction}

We will briefly describe two statistics for the measure of
dinucleotide over- and under-representation in sequences
\cite{Karlin,UV}, which can both be computed with \seqinr{}. We will
subsequently use them to answer the long-time controversial question
concerning the relationship between UV exposure and genomic content in
bacteria \cite{Singer,Bak}.

\subsection{The \textit{rho} statistic}

The $\rho$ statistic (\texttt{rho()}), presented in \cite{Karlin},
measures the over- and under-representation of two-letter words:

$$\rho(xy) = \frac{f_{xy}}{f_{x}\times f_{y}}$$

where $f_{xy}$ and $f_{x}$ are respectively the frequencies of
dinucleotide $xy$ and nucleotide $x$ in the studied sequence. The
underlying model of random generation considers dinucleotides to be
formed according to the specific frequencies of the two nucleotides
that compose it ($\rho_{xy} = 1$). Departure from this value
characterizes either over- or under-representation of dinucleotide
$xy$.


We expect the $\rho$ statistic of a randomly generated sequence to be
neither over- nor under-represented. Indeed, when we compute the
$\rho$ statistic on 11 random
sequences, we can fit a normal distribution which is centered on 1
(see Fig. \ref{rho})

\begin{figure}[H]
  \begin{center}
\begin{Schunk}
\begin{Sinput}
 rhoseq <- sapply(seq(n), function(x) {
     rho(sample(s2c("acgt"), 6000, rep = TRUE))
 })
 hist(rhoseq[di, ], freq = FALSE, xlab = "Rho statistic", main = paste("Distribution for dinucleotide", 
     toupper(labels(rhoseq)[[1]][di]), "on", n, "random sequences"), 
     las = 1, col = grey(0.8), border = grey(0.5))
 abline(v = 1, lwd = 2, lty = 3)
 min = min(rhoseq[di, ])
 max = max(rhoseq[di, ])
 x = seq(min, max, length.out = 1000)
 lines(x, dnorm(x, mean = mean(rhoseq[di, ]), sd = sd(rhoseq[di, 
     ])), lwd = 2, lty = 2)
\end{Sinput}
\end{Schunk}
\end{center}
   \caption{Distribution of the $\rho$ statistic computed on
     11 random sequences. The vertical
     dotted line is centered on 1. The dashed curve draws the fitted normal
     distribution.}
   \label{rho}
   \end{figure}

The downside of this statistic, is that the model against which we
compare the sequence under study is fixed. For several types of
sequences, dinucleotides are far from being formed by mere chance
(CDS, ...). In this case, the model used in the $\rho$ statistic
becomes trivial, and the over- or under-representations measured are
mainly due to the strong constraints acting on those sequences.

\subsection{The $z$-score statistic}

The $z$-score statistic (\texttt{zscore()}) is inspired by the
$\rho$ statistic, and is defined so that several different models can
be used for the determination of over- and under-representation
\cite{UV}. It allows for a finer measure of over- and
under-representation in sequences, according to the chosen model.

The $z$-score is defined as follows:

$$z_{score}=\frac{\rho_{xy}-E(\rho_{xy})}{\sqrt{Var(\rho_{xy})}}$$

where $E(\rho_{xy})$ and $Var(\rho_{xy})$ are the expected mean and
variance of $\rho_{xy}$ according to a given model that describes the
sequence.

This statistic follows the standard normal distribution, and can be
computed with several different models of random sequence generation
based on permutations from the original sequence (\texttt{modele}
argument). More details on those models can be obtained in the
documentation for the \texttt{zscore()} function, by simply typing
\texttt{?zscore}.

For instance, if we want to measure the over- and under-representation
of dinucleotides in CDS sequences, we can use the \texttt{codon}
model, which measures the over- and under-representations existing in
the studied sequence once codon usage bias has been erased. For
intergenic sequences, or sequences for which no good permutation model
can be established, we can use the \texttt{base} model. 

\subsection{Comparing statistics on a sequence}

Let's have a look at what these different statistics can show. First,
we will extract a CDS sequence of \textit{Escherichia coli}'s
chromosome from the Genome Reviews database. We will first make a
request to retrieve all available CDS from this bacteria:

\begin{Schunk}
\begin{Sinput}
 choosebank("greview")
 query("coli", "sp=escherichia coli et t=cds et no k=partial")
 names <- unlist(lapply(coli$req, getName))
 sequence <- getSequence(coli$req[[which(names == "U00096.PE448")]])
\end{Sinput}
\end{Schunk}

From the ??? sequences
annotated as CDS in the Genome Reviews database, let's choose one
coding sequence: say, for instance, the one with accession number \texttt{U00096.PE448}. We can see
that this CDS encodes a maltose O-acetyltransferase protein
(\verb+getAnnot(coli$req[[which(names=="U00096.PE448")]],30)+). We will now compare the three
following nonparametric statistics:

\begin{itemize}
\item the $\rho$ statistic,
\item the $z$-score statistic with \texttt{base} model,
\item and the $z$-score statistic with \texttt{codon} model.
\end{itemize}

The $z$-score statistic has been modified to incorporate an exact
analytical calculation of the \texttt{base} model where the old 
version (seqinR 1.1-1 and previous versions) incorporated an
approximation for large sequences. This has been possible with the
help of Sophie Schbath \cite{Schbath-thesis}, and the new version of
this calculation can be obtained with the argument \texttt{exact} 
set to \texttt{TRUE} (\texttt{FALSE} being the default).

\begin{figure}[H]
  \begin{center}
\begin{Schunk}
\begin{Sinput}
 rhocoli = rho(sequence)
 zcolibase = zscore(sequence, mod = "base", exact = TRUE)
 zcolicodon = zscore(sequence, mod = "codon")
 par(mfrow = c(1, 3))
 plot(rhocoli - 1, ylim = c(-0.5, 0.5), las = 1, ylab = "rho")
 plot(zcolibase, ylim = c(-2.5, 2.5), las = 1, ylab = "zscore with base model")
 plot(zcolicodon, ylim = c(-2.5, 2.5), las = 1, ylab = "zscore with codon model")
\end{Sinput}
\end{Schunk}
\end{center}
   \caption{Three different non-parametric statistics (from left to
   right: $\rho$, $zscore$ with \texttt{base} model, $zscore$ with
   \texttt{codon} model), computed on the same sequence from
   \textit{Escherichia coli}. In order to make the figures easily
   comparable, we substracted 1 to the \texttt{rho()} results, so that
   all 3 statistics are centered on 0. }
\label{dinuclcoli}
\end{figure}

The first two figures are almost identical: this is due to the way the
$z$-score statistic has been built. The statistic computed with
the \texttt{base} model is a reflection of the $\rho$ statistic. The
difference being that the $z$-score follows a standard normal
distribution, which makes easier the comparisons between the results
from the \texttt{base} model and the ones from the \texttt{codon}
model. The last figure ($z$-score with \texttt{codon} model), is
completely different: almost all over- and under-representations have
been erased. We can safely say that these over- and
under-representations were due to codon usage bias.

On this last figure, four dinucleotides stand out: CC and TT seem
rather under-represented, CT and TC rather over-represented. This
means that, in this sequence, codons ending with a given pyrimidine
tend to be more frequently followed by a codon starting with the other
pyrimidine than expected by chance. This is not a universal feature of
\textit{Escherichia coli}, and is probably due to the amino-acid
composition of this particular sequence. It seemed a funny example, as
the following part will also relate to pyrimidine dinucleotides.
However, what we see on this CDS from \textit{Escherichia coli} has
nothing to do with what follows...

\section{UV exposure and dinucleotide content}

In the beginning of the 1970's, two contradictory papers considered
the question of the impact of UV exposure on genomic content. Both
papers had strong arguments for either side, and the question remained
open until recently \cite{UV}.

\subsection{The expected impact of UV light on genomic content}

On this controversy, the known facts are: pyrimidine dinucleotides
(CC, TT, CT and TC) are the major DNA target for UV-light
\cite{Setlow}; the sensitivities of the four pyrimidine dinucleotides
to UV wavelengths differ and depend on the micro-organism
\cite{Setlow}:

\begin{table}[H]
\begin{center}
\begin{tabular}{rcccc}
  \hline
  & G+C content & CC (\%) & CT + TC (\%) & TT (\%)\\
  \hline
  \textit{Haemophilus influenzae} & 62 & 5 & 24 & 71 \\
  \hline
  \textit{Escherichia coli} & 50 & 7 & 34 & 59 \\
  \hline
  \textit{Micrococcus lysodeikticus} & 30 & 26 & 55 & 19 \\
  \hline
\end{tabular}
\caption{Proportion of dimers formed in the DNA of three bacteria
  after irradiation with 265 nm UV light. Table adapted from
  \cite{Setlow}.}
\label{sensitivity}
\end{center}
\end{table}

%% %% \begin{figure}
%% %%    \begin{minipage}[c]{.46\linewidth}
%% <<sensitivity1, eval = TRUE, echo = FALSE, fig = TRUE,eval=F>>=
%% pie(c(0.59,0.34,0.07),labels=c('T-T(59\%)','C-T(34\%)','C-C(7\%)'),col=c('lightblue','purple','royalblue'),main="Escherichia coli")
%% @
%% %%   \end{minipage} \hfill
%% %%    \begin{minipage}[c]{.46\linewidth}
%% <<sensitivity2, eval = TRUE, echo = FALSE, fig = TRUE,eval=F>>=
%% pie(c(0.71,0.24,0.05),labels=c('T-T(71\%)','C-T(24\%)','C-C(5\%)'),col=c('lightblue','purple','royalblue'),main="Haemophilus influenzae")
%% @ 
%% %%    \end{minipage} \hfill
%% %%    \begin{minipage}[c]{.46\linewidth}
%% <<sensitivity3, eval = TRUE, echo = FALSE, fig = TRUE,eval=F>>=
%% pie(c(0.19,0.55,0.26),labels=c('T-T(19\%)','C-T(55\%)','C-C(26\%)'),col=c('lightblue','purple','royalblue'),main="Micrococcus lysodeikticus")
%% @
%% %%    \end{minipage}
%% %% \end{figure}

The hypothesis presented by Singer and Ames \cite{Singer} is that
pyrimidine dinucleotides are avoided in light-exposed micro-organisms.
At the time, only G+C content is available, and -- based exclusively
on the sensitivity of the four pyrimidine dinucleotides in an
\textit{Escherichia coli} chromosome -- they hypothesize that a high
G+C will result in less pyrimidine target. Indeed, they find that
bacteria exposed to high levels of UV have higher G+C content than the
others. Bak \textit{et al.} \cite{Bak} strongly criticize their
methodology, but no clear cut answer is achieved.

In an \textit{Escherichia coli} chromosome, it is true that a sequence
with a high G+C content will contain few phototargets (see Fig.
\ref{GCcoli}).

\begin{figure}[H]
  \begin{center}
\end{center}
\caption{Density of phototargets, weighted by their frequency in the
  \textit{Escherichia coli} chromosome, and calculated for different
  G+C contents and for three kinds of random genomes. The weights are
  as follows: $0.59*f_{tt}+0.34*(f_{tc}+f_{ct})+0.07*f_{cc}$ (where
  $f_{xy}$ is the frequency of dinucleotide $xy$ in the specified
  genome). Three models of random genomes are analyzed. In the worst
  case (solid curve), the genome is the concatenation of a sequence of
  pyrimidines and a sequence of purines: all pyrimidines are involved
  in a pyrimidine dinucleotide. In the best case (dotted curve), the
  genome is an unbroken succession of pyrimidine-purine dinucleotides:
  no pyrimidine is involved in a pyrimidine dinucleotide. In the
  "random case" (dashed curve), the frequency of a pyrimidine
  dinucleotide is the result of chance ($f_{xy} = f_{x}\times
  f_{y}$).}
\label{GCcoli}
\end{figure}


In a \textit{Micrococcus lysodeikticus} sequence (see Fig.
\ref{GCmicro}), we can see that this is no longer true...


\begin{figure}[H]
   \begin{center}
\end{center}
\caption{Density of phototargets, weighted by their frequency in the
  \textit{Micrococcus lysodeikticus} chromosome, and calculated for
  different G+C contents and for three kinds of random genomes. The
  weights are as follows:
  $0.19*f_{tt}+0.55*(f_{tc}+f_{ct})+0.26*f_{cc}$. See previous figure
  for more details.}
\label{GCmicro}
\end{figure}

These two figures show that the density of phototargets depends on:

\begin{itemize}
\item the degree of aggregation of pyrimidine dinucleotides in the
sequence,
\item the sensitivities of the four pyrimidine dinucleotides.
  \end{itemize}

Instead of looking at G+C content, which is an indirect measure of the
impact of UV exposure on genomic content, let us look at pyrimidine
dinucleotide content.

Are CC, TT, CT and TC dinucleotides avoided in light-exposed bacteria?

\subsection{The measured impact of UV light on genomic content}

On all available genomes (as retrieved from Genome Reviews database on
June 16, 2005), we have computed the mean of the $z$-score with
the \texttt{base} model on all intergenic sequences, and the mean of
the $z$-score with the \texttt{codon} model on all CDS:

\begin{Schunk}
\begin{Sinput}
 data(dinucl)
\end{Sinput}
\end{Schunk}

The results show that there is no systematic under-representation of
none of the four pyrimidine dinucleotides (see Fig.
\ref{dinucl-graphes}).

\begin{figure}[H]
   \begin{center}
\end{center}
   \caption{Plot of the mean $zscore$ statistics for
   \textbf{intergenic sequences} (x-axis) and for \textbf{coding
   sequences} (y-axis), for each of the four pyrimidine dinucleotides.
   On each plot, a dot corresponds to the mean of these two statistics
   in a given prokaryote chromosome. The null x and y axis (dotted
   lines), and the 5\% limits of significance for the standard normal
   distribution (dashed lines) are plotted as benchmarks. It should be
   noted that the variability within one chromosome is sometimes as
   great as that between different chromosomes.}
\label{dinucl-graphes}
\end{figure}

However, we have little or no information on the exposure of this
bacteria to UV light. In order to fully answer this question, let's do
another analysis and look at \textit{Prochlorococcus marinus} genome.

\textit{Prochlorococcus marinus} seems to make an ideal model for
investigating this hypothesis. Three completely sequenced strains are
available in the Genome reviews database: two of these strains are
adpated to living at a depth of more than 120 meters (accession
numbers AE017126 and BX548175), and the other one at a depth of 5
meters (accession number BX548174).

Living at a depth of 5 meters, or at a depth of more than a 120 meters
is totally different in terms of UV exposure: the residual intensity
of 290 nm irradiation (UVb) in pure water can be estimated to 56\% of
its original intensity at 5 m depth and to less than 0.0001\% at more
than 120 m depth. For this reason, two of the \textit{Prochlorococcus
marinus} strains can be considered to be adapted to low levels of UV
exposure, and the other one to much higher levels. Is pyrimidine
dinucleotide content different in these three strains? And is it
linked to their UV exposure?

We have computed the $z$-score with the \texttt{codon} model on
all CDS from each of these three strains (as retrieved from Genome
Reviews database on June 16, 2005):

\begin{Schunk}
\begin{Sinput}
 data(prochlo)
\end{Sinput}
\end{Schunk}


\begin{figure}[H]
   \begin{center}
\end{center}
   \caption{Each figure shows the distributions of the $zscore$ in all
     \textbf{coding sequences} corresponding to each of the three
     strains of \textit{Prochlorococcus marinus}. In each figure, the
     distribution for the MED4 (a high-light adapted strain) is shown
     as a solid line; the distribution for the SS120 (a low-light
     adapted strain) is shown as a dashed line, and the distribution
     for the MIT 9313 (a low-light adapted strain) is shown as a
     dotted line. The 5\% limits of significance for the standard
     normal distribution (dashed vertical lines) are plotted as
     benchmarks.}
\label{prochlo}
\end{figure}

Figure \ref{prochlo} shows that there is no difference between the
relative abundances of pyrimidine dinucleotides in these three
strains. We can say that pyrimidine dinucleotides are not avoided, and
that the hypothesis by Singer and Ames \cite{Singer} no longer stands
\cite{UV}.


\section*{Session Informations}

This part was compiled under the following \Rlogo{}~environment:

\begin{itemize}
  \item R version 2.6.0 (2007-10-03), \verb|i386-apple-darwin8.10.1|
  \item Locale: \verb|C|
  \item Base packages: base, datasets, grDevices, graphics, methods,
    stats, utils
  \item Other packages: MASS~7.2-36, ade4~1.4-4, ape~2.0-1,
    gee~4.13-13, lattice~0.16-5, nlme~3.1-85, quadprog~1.4-11,
    seqinr~1.1-3, tseries~0.10-11, xtable~1.5-1, zoo~1.4-0
  \item Loaded via a namespace (and not attached): grid~2.6.0,
    rcompgen~0.1-15
\end{itemize}
There were two compilation steps:

\begin{itemize}
  \item \Rlogo{} compilation time was: Sun Nov  4 16:47:06 2007
  \item \LaTeX{} compilation time was: \today
\end{itemize}


% END - DO NOT REMOVE THIS LINE

%%%%%%%%%%%%  BIBLIOGRAPHY %%%%%%%%%%%%%%%%%
\clearpage
\addcontentsline{toc}{section}{References}
\bibliographystyle{plain}
\bibliography{../config/book}
\end{document}
